%% 封面、中英文摘要、目录（自动生成，学生不要改这个文件）
%% 需要改题目、姓名、摘要时，改 大论文模板.qmd 顶部 YAML 区
%% 封面版式在 ../cover/cover.tex，与开题报告模板共用；这里只传文字
%% 页眉第二行“题目：”后自动填入论文题目
\providecommand{\zjgsuTitle}{}
\renewcommand{\zjgsuTitle}{$title$}
%% ctexart 在 \begin{document} 时会重置字号，这里再固定一次正文小四、20 磅
\renewcommand{\normalsize}{\fontsize{12pt}{20pt}\selectfont}
\fontsize{12pt}{20pt}\selectfont

%% 共用封面（2026-11-10）
%% 大论文模板与开题报告模板共用这一份：改版式只改这里，两边同时生效。
%% 调用方在各自 before-body.tex 里先定义下面几个宏，再调用 \zjgsuCover：
%%   \zjgsuCoverHeading  标题行文字，如“本科毕业论文（设计）正文”
%%   \zjgsuCoverTitleCN  中文题目（不含“题目：”三字）
%%   \zjgsuCoverTitleEN  英文题目；留空则不出这一行
%%   \zjgsuCoverRows     信息栏，每行写成“标　签：&值\cr”
%%   \zjgsuCoverDate     封面日期
%% 校徽、校名取自教务处 doc 原件 word/media；实测 校徽 3.65×3.41cm、校名 10.61×2.01cm

\providecommand{\zjgsuCoverHeading}{}
\providecommand{\zjgsuCoverTitleCN}{}
\providecommand{\zjgsuCoverTitleEN}{}
\providecommand{\zjgsuCoverRows}{}
\providecommand{\zjgsuCoverDate}{}
%% 信息栏的上留白、下留白、行距：默认值取自大论文 6 行时的成品，
%% 开题模板沿用默认即可；行数不同的调用方自行覆盖这三项，把字段区重新均分。
\providecommand{\zjgsuCoverRowsAbove}{0.7cm}
\providecommand{\zjgsuCoverRowsBelow}{1.1cm}
\providecommand{\zjgsuCoverRowSkip}{35pt}

\def\zjgsuCover{%
  \thispagestyle{empty}%
  \begin{center}%
  \vspace*{0.3cm}%
  \includegraphics[width=3.65cm]{../cover/zjgsu-logo.jpg}\par
  \vspace{0.25cm}%
  \includegraphics[width=10.6cm]{../cover/zjgsu-name.png}\par
  \vspace{1.4cm}%
  {\fontsize{26pt}{34pt}\selectfont\bfseries \zjgsuCoverHeading\par}%
  \vspace{1.0cm}%
  %% 《要求》封面：题目/英文题目三号加粗，行距固定 35 磅
  {\fontsize{16pt}{35pt}\selectfont\bfseries 题目：\zjgsuCoverTitleCN\par}%
  \ifx\zjgsuCoverTitleEN\empty\else
    {\fontsize{16pt}{35pt}\selectfont\bfseries \zjgsuCoverTitleEN\par}%
  \fi
  \vspace{\zjgsuCoverRowsAbove}%
  %% 信息栏 16pt 加粗，行距由 \zjgsuCoverRowSkip 给；整块（标签列＋值列）以版心为中心居中
  {\fontsize{16pt}{\zjgsuCoverRowSkip}\selectfont\bfseries
  \setbox0=\vbox{\tabskip=0pt\halign{\hfil##\quad&##\hfil\cr \zjgsuCoverRows}}\par
  \begingroup\centering\noindent\box0\par\endgroup}%
  \vspace{\zjgsuCoverRowsBelow}%
  {\fontsize{16pt}{20pt}\selectfont\bfseries \zjgsuCoverDate\par}%
  \end{center}%
  \newpage
}

\def\zjgsuCoverHeading{本科毕业论文（设计）正文}
\def\zjgsuCoverTitleCN{$title$}
\def\zjgsuCoverTitleEN{$subtitle$}
%% 《正文模板》段 18–20：学院/专业/班级/学号/姓名/指导教师
%% 2026-11-07 主人明令字段值「连同标签一起居中」——整块（标签列＋值列）
%% 以版心为中心居中，标签仍右对齐、值仍左对齐，列内不错位。
%% 2026-11-10：按主人指令**去掉“学院”一行**（学校原件封面是有学院栏的，此处为有意去掉，
%% 不是漏改，别再加回去）。学院值仍留在 .qmd 的 college 键里，只是不再渲染。
%% 行数由 6 变 5 后重新均分字段区：省下的一行 35 磅，用来把字段区上下留白补成相等
%% （原来上 20.2、下 35.2 磅），其余摊到 4 个行间；实测 上下留白 34.7、行距 40.0 磅，
%% 字段区整体占位与日期位置仍与 6 行时重合。
\def\zjgsuCoverRowsAbove{34.40pt}
\def\zjgsuCoverRowsBelow{30.85pt}
\def\zjgsuCoverRowSkip{40pt}
\def\zjgsuCoverRows{专　　业：&$major$\cr
班　　级：&$class$\cr
学　　号：&$studentid$\cr
学生姓名：&$author$\cr
指导教师：&$advisor$\cr}
\def\zjgsuCoverDate{$yearmonth$}
\zjgsuCover

%% ---------- 中文摘要 ----------
%% 范例：封面不编页码；摘要、目录用阿拉伯数字 1、2、3；正文重新从 1 开始
\pagenumbering{arabic}
\setcounter{page}{1}
%% 范例段 23：中文题目三号黑体加粗居中，上下各空一行；
%% 标题→摘要实测 57.88 磅（模板原先 43.48，补 14.4 磅）
\vspace*{20pt}
\begin{center}
{\heiti\fontsize{16pt}{24pt}\selectfont\bfseries $title$\par}
\end{center}
\vspace{34.4pt}
{\xiaosi\noindent\hspace*{2em}\textbf{摘\quad 要：}$abstract$\par}

%% 范例摘要末行→关键词 20.00 磅，不加额外空行
{\xiaosi\noindent\hspace*{2em}\textbf{关键词：}$keywords$\par}
%% 2026-11-07 主人明令：中文摘要之后不强制换页。范例第 2 页关键词（439.79）
%% 与英文题目（500.54）同页接排，墨迹间距 60.75 磅；此处按同一间距接排。
\vspace*{14.70pt}

%% ---------- 英文摘要 ----------
\begin{center}
%% 范例段 28：英文题目居中、16pt 加粗，w:line="360" w:lineRule="auto"（1.5 倍行距），
%% LibreOffice 渲染为行距 46.80 磅（实测 500.54→547.34）
{\fontsize{16pt}{46.8pt}\selectfont\bfseries $subtitle$\par}
\end{center}
%% 范例第 2 页：英文题目末行 y=547.34 → Abstract 首行 y=622.36，75.02pt（实测）
\vspace{52.34pt}
%% 范例段 30：Abstract 首行缩进 2 字符 = 24.0pt（正文模板 306070 EMU=24.1pt）
{\xiaosi\noindent\hspace*{2em}\textbf{Abstract:} $abstract-en$\par}

{\xiaosi\noindent\hspace*{2em}\textbf{Keywords:} $keywords-en$\par}
\newpage

%% ---------- 目录 ----------
\renewcommand{\contentsname}{正文目录}
\vspace*{-5.6pt}
\tableofcontents
\newpage

%% ---------- 正文从 1 开始编页码 ----------
\pagenumbering{arabic}
\setcounter{page}{1}
\vspace*{-15.3pt}
